\documentclass[11pt]{article}
\usepackage[margin=1in]{geometry}
\usepackage{amsmath,amssymb,bm,braket}
\usepackage{hyperref}
\begin{document}
\paragraph*{Step-by-Step Derivation}

A vacuum expectation value of a Verlinde line equals its quantum dimension,

\[
\lambda_a=\langle L_a\rangle
=\frac{S_{a0}}{S_{00}},
\qquad
\lambda_0=1.
\]

This is the standard vacuum normalization for topological defect lines in a unitary CFT. \href{\detokenize{https://www.mdpi.com/1099-4300/26/12/1064}}{Entropic Order Parameters for Categorical Symmetries in 2D CFT}

The Moore--Read edge theory is the product of an Ising neutral sector and an Abelian charged boson sector. Its electron is the product of the Ising fermion and a charged vertex operator. \href{\detokenize{https://harvest.aps.org/v2/journals/articles/10.1103/RevModPhys.80.1083/fulltext}}{Nayak et al., *Non-Abelian Anyons and Topological Quantum Computation*}

For the Ising, equivalently $SU(2)_2$, sector,

\[
S^{\mathrm{Ising}}
=
\frac12
\begin{pmatrix}
1&\sqrt2&1\\
\sqrt2&0&-\sqrt2\\
1&-\sqrt2&1
\end{pmatrix},
\]

in the ordering $j=0,\frac12,1$. Therefore,

\[
d_j=\frac{S_{j0}}{S_{00}}
=
\frac{\sin\!\left[\frac{\pi}{4}(2j+1)\right]}
{\sin(\pi/4)}
=
\begin{cases}
1,&j=0,\\[2mm]
\sqrt2,&j=\frac12,\\[2mm]
1,&j=1.
\end{cases}
\]

The general $SU(2)_k$ expression for these vacuum Verlinde-line expectation values is given explicitly in \href{\detokenize{https://scoap3-prod-backend.s3.cern.ch/media/files/67899/10.1007/JHEP01%282022%29175_a.pdf}}{JHEP 01 (2022) 175}.

Every line in the Abelian $U(1)$ charge sector has quantum dimension one, so the answer is independent of $n_L,n_R$. \href{\detokenize{https://s3.cern.ch/inspire-prod-files-e/e3ba3709c35f55bab8ac4ed87349533e}}{JHEP 10 (2015) 092}

Since the left- and right-moving modular matrices factorize,

\[
\lambda_{(j_L,n_L,j_R,n_R)}
=d_{j_L}d_{j_R}
=2^{\frac12\delta_{j_L,1/2}+\frac12\delta_{j_R,1/2}}.
\]

For the physical Moore--Read spectrum, locality with respect to each electron requires

\[
2j_L+n_L\equiv0\pmod 2,
\qquad
2j_R+n_R\equiv0\pmod 2.
\]

Thus each chiral edge has the six sectors

\[
\mathcal A=
\left\{
(0,0),(0,2),(1,0),(1,2)
\right\},
\qquad
\mathcal N=
\left\{
\left(\frac12,1\right),
\left(\frac12,3\right)
\right\},
\]

where sectors in $\mathcal A$ have dimension $1$, while those in $\mathcal N$ have dimension $\sqrt2$. The electron lines consequently have expectation value $1$, as required.

\textbf{Final Answer:}

\[
\boxed{
\begin{aligned}
\Big\{&(j_L,n_L,j_R,n_R,1):
 (j_L,n_L)\in\mathcal A,\ (j_R,n_R)\in\mathcal A\Big\}\\
{}\cup\Big\{&(j_L,n_L,j_R,n_R,\sqrt2):
 (j_L,n_L)\in\mathcal A,\ (j_R,n_R)\in\mathcal N
 \ \text{or vice versa}\Big\}\\
{}\cup\Big\{&(j_L,n_L,j_R,n_R,2):
 (j_L,n_L)\in\mathcal N,\ (j_R,n_R)\in\mathcal N\Big\},
\end{aligned}}
\]

with

\[
\mathcal A=\{(0,0),(0,2),(1,0),(1,2)\},
\qquad
\mathcal N=\left\{\left(\frac12,1\right),\left(\frac12,3\right)\right\}.
\]

This gives $16$ tuples with $\lambda=1$, $16$ with $\lambda=\sqrt2$, and $4$ with $\lambda=2$.
\end{document}
